\chapter{Introduction}
\label{ch:introduction}

\section{Motivation}
The growing computational power of quantum computers poses a significant threat to current classical cryptosystems. For example, Shor's algorithm can efficiently solve the factoring problem, while Grover's search algorithm can effectively reduce the security level of symmetric-key cryptography by shortening its effective key length.

On the other hand, quantum computers also open new directions for the design of cryptographic and cryptanalytic techniques, allowing computation to be reconsidered from a physical perspective, a core principle underlying quantum computation and quantum information. In particular, quantum cryptography, as its name suggests, exploits cryptographic mechanisms rooted in quantum circuits. A key advantage of leveraging quantum mechanisms is that they provide information-theoretic secrecy guaranteed by the uncertainty principle.

To improve the efficiency of quantum circuits, it is essential to develop more compact implementations of their underlying constructions. In quantum computing, many applications rely heavily on linear reversible circuits; therefore, deploying efficient methods to reduce computational cost is key to enabling more practical algorithms.

Currently, most symmetric cryptosystems are constructed using a significant number of CNOT gates. Drawing inspiration from the field of Very Large Scale Integration (VLSI), we aim to reduce the cost of quantum algorithms by focusing on circuit depth as our primary metric, targeting linear scaling.

From a cryptographic engineering standpoint, asymptotic security claims such as
Grover's $O(2^{n/2})$ query complexity are necessary but insufficient for
practical risk assessment. Standardization bodies and practitioners tasked with
selecting, deprecating, or migrating cryptographic primitives require concrete,
comparable cost figures rather than asymptotic bounds alone, since two
primitives with identical asymptotic security can differ substantially in the
practical effort required to attack them once implementation-level
details are accounted for. This distinction underlies, for example, NIST's
post-quantum cryptography standardization process, which defines its five
security categories not purely in terms of asymptotic complexity but by
benchmarking a candidate scheme's concrete attack cost against the estimated
classical and quantum resources required to break AES-128, AES-192, and
AES-256~\cite{chen2016report}. Under this framework, refining the resource
estimate of a Grover-based attack on a given cipher as this thesis does for
KLEIN is not a peripheral exercise but a direct input to how that cipher
would be categorized and evaluated for deployment in a post-quantum threat
model.

This concern is particularly salient for lightweight ciphers such as KLEIN,
which are designed for deployment on resource-constrained devices: RFID
tags, sensor nodes, and other embedded Internet-of-Things (IoT) endpoints,
that are frequently expected to remain in service for a decade or longer after
deployment. Such devices are rarely field-upgradable, making cryptographic
agility difficult or impossible to retrofit after the fact. Consequently, the
security margin of a lightweight cipher against a future large-scale quantum
adversary must be assessed today, using the most accurate available
resource estimates, rather than revisited once quantum hardware matures. This
is further compounded by the ``harvest now, decrypt later'' threat model, in
which an adversary records encrypted or key-exchanged traffic in the present
with the intent of decrypting it retroactively once sufficiently capable
quantum hardware becomes available; for long-lived embedded deployments using
static or infrequently rotated keys, this threat model applies directly to
symmetric-key primitives such as KLEIN, not only to public-key schemes.

In this context, the present work contributes a tighter, more accurate lower
bound on the practical quantum resources required to mount a Grover-based
exhaustive key search against KLEIN, obtained through a genuine circuit-level
optimization rather than a change in assumptions or methodology. As discussed
in Section~\ref{subsec:grover-implication}, this refinement operates entirely
within the concrete resource-cost axis of security assessment, leaving the
asymptotic query complexity, and therefore the cipher's nominal security
level in bits, unchanged. Nonetheless, because resource-based cost estimates
of this kind directly inform standardization decisions, migration timelines,
and parameter selection for post-quantum cryptographic deployment, providing
increasingly accurate estimates constitutes a meaningful contribution to the
broader information security objective of ensuring that deployed
cryptographic primitives possess an accurately characterized, rather than
merely asymptotically bounded, security margin against quantum adversaries.

In this work, we propose two enhanced depth-oriented greedy algorithms that
employ heuristic methods to achieve the shallowest possible circuit depth,
while producing a functionally equivalent circuit with a modest reduction in
CNOT gate count. Our approach not only guarantees a speedup in the execution
of quantum cryptographic algorithms but also reduces the impact of quantum
noise and hardware constraints on computational results.

\section{Problem Statement}
In this study, we address the challenge of finding the optimal (i.e., shallowest) depth for constructing the matrix components of symmetric ciphers in gate-based quantum circuits. All matrices are defined over $\mathbb{F}_{2}$. We subsequently verify that our results are applicable to real-world block cipher schemes.

\section{Contributions}
Our research focuses on optimizing the quantum circuits of block ciphers. Specifically, components that provide diffusion and confusion properties for symmetric cryptography. We survey the latest quantum circuit implementations of lightweight cryptosystems, and demonstrate that our improved AES MixColumns (Rijndael) circuit significantly reduces the number of CNOT gates required, a result that also extends to the KLEIN family.

All of source codes and results of this paper are published in: https://github.com/\\hshuang99/Depth-Oriented-Greedy-Algorithms-for-Optimizing-Linear-Reversible-\\Quantum-Circuits-of-Block-Ciphers

\section{Organization}
The remainder of this paper is organized as follows. Chapter~\ref{ch:preliminaries} introduces the preliminary concepts and notation used throughout this work. Chapter~\ref{ch:existing-approaches} reviews existing approaches to constructing optimal quantum circuits. Chapter~\ref{ch:our-algorithms} presents our proposed greedy algorithm. Chapter~\ref{ch:klein} introduces the lightweight cipher family KLEIN and demonstrates that our optimized circuits are applicable to it. Chapter~\ref{ch:experimental-results} presents our experimental results. Finally, Chapter~\ref{ch:discussion} and \ref{ch:conclusion} provide discussion and conclusions, respectively.